\documentclass[12pt,letterpaper]{article}

% --- Paquetes Requeridos ---
\usepackage[utf8]{inputenc}
\usepackage[spanish,es-tabla]{babel}
\usepackage[margin=2.5cm]{geometry}
\usepackage{graphicx}
\usepackage{float}
\usepackage{titlesec}
\usepackage{setspace}
\usepackage{hyperref}
\usepackage{booktabs}
\usepackage{caption}
\usepackage{enumitem}
\usepackage{listings}
\usepackage{xcolor}

% --- Configuración de Párrafo ---
\onehalfspacing
\setlength{\parindent}{0pt}
\setlength{\parskip}{1em}

% --- Rutas de búsqueda de imágenes ---
% Permite compilar tanto desde Documentos/Entregables/ como con todas las
% imágenes en una sola carpeta plana.
\graphicspath{{./}{../Evidencia/}{Evidencia/}{imagenes/}}

% --- Configuración de listings ---
\lstset{
    basicstyle=\ttfamily\footnotesize,
    breaklines=true,
    frame=single,
    captionpos=b,
    showstringspaces=false
}

% --- Inicio del Documento ---
\begin{document}

% --- PORTADA ---
\begin{titlepage}
    \centering

    {\Large \textbf{UNIVERSIDAD SANTO TOMÁS}} \\[1cm]

    {\Large \textbf{Sistema colaborativo de robots móviles para asistencia de orientación en entornos interiores con múltiples pisos}} \\[2cm]

    {\large \textbf{Realizado por:}} \\
    {\large Santiago Hernández Ávila} \\
    {\large Jonny Alejandro Mejía León} \\[1cm]

    \begin{figure}[H]
        \centering
        \includegraphics[width=0.6\textwidth]{Logos Uni.png}
        \label{Logo Universidad Santo Tomas y facultad de ingenieria electronica}
    \end{figure}

    \vfill

    {\large \textbf{Informe de Avance --- Semana 23}} \\
    {\large Fase 5: integración del sistema y realización de pruebas} \\
    {\large Cierre de la implementación, llegada del segundo vehículo y medición de los pasillos reales} \\[1.5cm]

    \textbf{Dirigido por:} \\
    Ing. Armando Mateus Rojas, Msc. \\
    Ing. Nestor Ivan Ospina, Msc. \\
    Ing. Oscar Mauricio Gélvez Lizarazo, Msc. \\[1.0cm]

    \vfill

    {\large GED (Grupo de Estudio y Desarrollo en Robótica)} \\
    {\large Facultad de Ingeniería Electrónica} \\
    {\large División de Ingenierías} \\
    {\large Bogotá, septiembre de 2026}
\end{titlepage}

% --- CONTENIDO ---

\section{Introducción}

Este informe presenta el avance de la semana 23 del cronograma (14 al 20 de septiembre de 2026), dentro de la fase 5, integración del sistema y realización de pruebas.

En esta semana se cerró la implementación. El 18 de septiembre se congeló el código con la etiqueta \texttt{v0.4-implementacion-congelada} (hito H7 del cronograma), y desde entonces no se añade funcionalidad nueva. En esa fecha estaban completos los objetivos específicos 1 (arquitectura funcional) y 3 (interfaz humano--robot). Al objetivo 4 (evaluación experimental) solo le faltaba la demostración física, que depende de la plataforma del objetivo 2.

El 14 de septiembre llegó el segundo vehículo, y ese mismo día se verificó la comunicación entre los dos (requisito RF-15), con lo que el objetivo 2 avanzó por primera vez en cuatro semanas (sección~\ref{sec:red}). Se verificó también la cancelación de misiones (RF-29), el último requisito pendiente del objetivo 3 (sección~\ref{sec:rf29}). Por último, se midió la información de avance en los pasillos reales de los dos pisos. En ninguno alcanza para que el vehículo construya el mapa mientras recorre el pasillo, así que la única ruta crítica que queda es publicar la odometría del vehículo (sección~\ref{sec:pisos}).

De las tres partes del criterio de cierre de la semana se cumplieron dos. La corrida física completa no se realizó (sección~\ref{sec:cronograma}).

\section{Objetivos de la semana}

El cronograma fija para esta semana el siguiente criterio de cierre:

\begin{quote}
\textit{El sistema corre de extremo a extremo en simulación y se ejecuta al menos una corrida física completa. Repositorio etiquetado.}
\end{quote}

Las actividades planificadas eran cuatro: verificación del funcionamiento general y corrección de fallas, ensayo del protocolo con los vehículos reales, congelación del código y emisión del informe de la semana 22. El criterio se cumplió en dos de sus tres partes; la corrida física completa no se realizó, y en la sección~\ref{sec:cronograma} se declara como no cumplida.

\section{Cancelación de misiones y verificación de RF-29}
\label{sec:rf29}

La interfaz tenía un botón para cancelar la misión desde el 2 de septiembre, pero no tenía efecto. En ROS~2, un servidor de acciones rechaza por defecto las solicitudes de cancelación si no se le programa una función que las atienda, y el coordinador no la tenía. Jonny Mejía corrigió el defecto el fin de semana anterior. Además, programó que una misión cancelada devuelva el vehículo a la escalera de su piso, que es el punto de transferencia, en lugar de dejarlo detenido en mitad del pasillo.

Con esa corrección se redactó el requisito RF-29 (el usuario cancela una misión en curso y el vehículo regresa al punto de transferencia de su piso) y se verificó con dos corridas (Tabla~\ref{tab:rf29}).

\begin{table}[H]
    \centering
    \small
    \caption{Corridas de verificación de RF-29.}
    \label{tab:rf29}
    \begin{tabular}{@{}p{3.6cm}p{10.4cm}@{}}
        \toprule
        \textbf{Corrida} & \textbf{Resultado} \\ \midrule
        Con cancelación & El registro marca el estado \texttt{CANCELANDO} y la misión cierra como \texttt{FALLIDA}. El vehículo termina a 0,121~m de la escalera; el criterio es 0,25~m. \\ \addlinespace
        De control, sin cancelación & Ninguna marca de cancelación; la misión cierra como \texttt{COMPLETADA}. \\ \bottomrule
    \end{tabular}
\end{table}

En la corrida con cancelación, la orden de cancelar se envía cuando la odometría indica que el vehículo se ha alejado más de 8~m, y no después de un tiempo fijo. El vehículo parte a 1,4~m de la escalera, así que una cancelación temprana cumpliría el criterio de regreso sin que el vehículo se hubiera desplazado.

La ejecución de la prueba mostró un error en el criterio de verificación. El criterio pedía comprobar que el registro de la misión incluyera el motivo ``Cancelada por el usuario'', pero ese texto se envía en la respuesta de la acción de ROS~2 y no queda guardado en el registro. El criterio se corrigió antes de dar el requisito por verificado.

\section{Llegada del segundo vehículo y verificación de RF-15}
\label{sec:red}

El 14 de septiembre el codirector entregó el segundo vehículo (\texttt{amss-jgm9}), que estaba en intervención técnica desde agosto (riesgo R11). Ese mismo día se midió la comunicación entre los dos vehículos para el requisito RF-15, con los límites fijados antes de la medición (Tabla~\ref{tab:rf15}).

\begin{table}[H]
    \centering
    \small
    \caption{Medición de RF-15 entre los dos vehículos, 14 de septiembre.}
    \label{tab:rf15}
    \begin{tabular}{@{}p{5.6cm}p{4.4cm}p{3.6cm}@{}}
        \toprule
        \textbf{Magnitud} & \textbf{Medido} & \textbf{Límite} \\ \midrule
        Mensajes recibidos & 600 de 600 (0~\% de pérdida) & 0~\% de pérdida \\ \addlinespace
        Tiempo de ida y vuelta, mediana & 7,61~ms & 100~ms \\ \addlinespace
        Tiempo de ida y vuelta, percentil 95 & 20,76~ms & 250~ms \\ \bottomrule
    \end{tabular}
\end{table}

La mediana queda 13 veces por debajo del límite, así que la red no restringe la publicación del estado del sistema, que se hace a 2~Hz. Con esta medición se verificó RF-15, se cerró el riesgo R11 después de 31 días abierto y el avance del objetivo 2 pasó del 55 al 65~\%.

Durante la prueba se encontraron dos problemas que el guion no preveía. El cortafuegos de los vehículos bloqueaba el tráfico entre ellos en los dos sentidos: la prueba de conectividad (\texttt{ping}) respondía, pero ninguno veía los nodos del otro. Por eso RF-15 no pudo verificarse durante un mes, aunque solo una de las dos causas era la falta del segundo vehículo. Además, el servidor DHCP asigna una dirección IP distinta en cada sesión, así que desde entonces cada vehículo se identifica por su nombre de equipo y su dirección MAC.

\section{Información de avance en los pasillos reales}
\label{sec:pisos}

El vehículo no tiene codificadores en las ruedas, así que su odometría se estima con rf2o, un paquete que compara barridos consecutivos del LiDAR. Para que rf2o detecte el avance, parte de los rayos deben incidir sobre superficies orientadas hacia la dirección de marcha, porque solo esos rayos cambian de longitud cuando el vehículo avanza. En un pasillo recto de paredes uniformes casi todos los rayos caen sobre las paredes laterales, y el avance no se registra.

La información de avance es la fracción de rayos de un barrido que inciden sobre superficies orientadas hacia la marcha, es decir, cuya normal forma menos de 45$^\circ$ con la dirección de avance. Se calcula con la herramienta \texttt{medir\_informacion\_avance.py}. Como referencia, en una caja cerrada simulada, cuyo mapa se construyó correctamente, vale 13,8~\%, y en el pasillo simulado, cuyo mapa falló, 6,8~\%. En los pasillos reales se midió lo que muestra la Tabla~\ref{tab:pisos}.

\begin{table}[H]
    \centering
    \small
    \caption{Información de avance en los pasillos reales.}
    \label{tab:pisos}
    \begin{tabular}{@{}p{6.0cm}p{6.4cm}@{}}
        \toprule
        \textbf{Entorno} & \textbf{Información de avance} \\ \midrule
        Pasillo real del piso 1 & 5,1~\% (mediana de 700 barridos) \\ \addlinespace
        Pasillo real del piso 2 & 5,9~\% (mediana de 849 barridos) \\ \bottomrule
    \end{tabular}
\end{table}

Antes de medir el piso 2 se registró una predicción de entre 8 y 14~\%, porque su pasillo tiene el doble de discontinuidades en las paredes que el del piso 1 (cuatro frente a dos). La medición no la confirmó. La Figura~\ref{fig:decaimiento} muestra la razón: la información que aporta el hall cae a partir de unos 6~m de distancia, y no a los 10 a 12~m del alcance del sensor, como se había supuesto.

\begin{figure}[H]
    \centering
    \includegraphics[width=0.95\textwidth]{S23_decaimiento_informacion_avance.png}
    \caption{Información de avance en el pasillo del piso 2 según la distancia al hall: 17,5~\% a 2,2~m y 5,9~\% a partir de 6,2~m. El umbral coincide con el medido en la semana 22 con otro método.}
    \label{fig:decaimiento}
\end{figure}

Este resultado no impide navegar el edificio. La campaña de evaluación en simulación, sobre el modelo de esos mismos pasillos, alcanzó un 86,7~\% de éxito con el mapa conocido de antemano. Lo que el pasillo no permite es que el vehículo construya el mapa mientras lo recorre.

Al revisar una corrida del piso 2 que no pudo usarse se encontró, además, que el LiDAR está montado girado 180$^\circ$, con lo que su sector ciego de 60$^\circ$ queda hacia delante.

\section{Respuesta de tracción del vehículo (RF-14)}
\label{sec:rf14}

El 17 de septiembre se midió por primera vez, sobre el vehículo, cómo se convierten las órdenes de velocidad de ROS~2 (tópico \texttt{/cmd\_vel}) en potencia del motor. La medición corresponde al requisito RF-14, según el cual cada vehículo se comanda desde ROS~2 sin pasar por la interfaz web del fabricante. Se encontraron dos defectos. El primero es una banda muerta: toda orden inferior a 0,40~m/s produce potencia cero. El segundo es que el nodo que convierte las órdenes aplica después un reescalado que reduce los tres niveles de potencia a 0,4247, 0,6242 y 0,7341 (en una escala de 0 a 1). En la misma sesión se recalibró la dirección, que nunca se había centrado.

Durante la sesión ocurrió un incidente de seguridad: una sola orden de tracción movió los dos vehículos. Los dos están en la misma red de ROS~2 y el tópico de las órdenes de servo no lleva espacio de nombres (un prefijo como \texttt{/robot1} que distinga a cada vehículo), así que ambos recibieron la orden. Es el pendiente de espacios de nombres del objetivo 2, ahora con consecuencias físicas. Desde entonces, en las pruebas de campo solo se enciende un vehículo.

\section{Congelación del código y plan del trabajo restante}
\label{sec:congelacion}

El 18 de septiembre se congeló el código con la etiqueta \texttt{v0.4-implementacion-congelada} y se publicó \texttt{MAPA\_TRABAJO\_RESTANTE.md}, que reúne en un solo documento el trabajo pendiente hasta la sustentación. Según ese análisis, el proyecto tiene una sola ruta crítica: publicar en el vehículo real la transformada \texttt{odom} $\rightarrow$ \texttt{base\_link}, es decir, la estimación de la posición del vehículo respecto a su punto de partida. Los seis requisitos que siguen abiertos dependen todos de la plataforma física.

\section{Otros avances}
\label{sec:otros}

El 16 de septiembre se cerró el riesgo R14 (dos criterios del protocolo experimental no podían dar un resultado negativo) con sus siete acciones. Dos de ellas se ejecutaron de forma distinta a como estaban enunciadas, con la razón documentada. Una de las acciones preveía regenerar los registros de las misiones a partir de las grabaciones de datos de ROS (\textit{rosbag}). Se comprobó que hacerlo reescribía la información de procedencia de los 44 registros, así que no se regeneró ninguno. La misma comprobación mostró que la regeneración es reproducible: los 44 registros se obtienen de nuevo a partir de las grabaciones sin que cambie ningún veredicto.

Desde esta semana, el repositorio guarda la fuente LaTeX de cada entregable en lugar del PDF, que se compila en Overleaf.

Se documentó en \texttt{GUIA\_ARRANQUE.md} la secuencia completa de arranque del sistema, que hasta entonces solo conocían los dos integrantes del equipo, para que otra persona pueda ejecutarla.

El 19 de septiembre se escribió el capítulo de resultados de la campaña en simulación, previsto para la semana 24. Los cuatro fallos de la campaña tienen la misma forma. Nav2 informa que llegó a la meta, pero la posición real del vehículo en el simulador queda entre 0,269 y 0,345~m de ella, por fuera de la tolerancia de 0,25~m. El capítulo atribuye el error a la estimación de la posición: el presupuesto de error del protocolo le asigna 88~mm, y en esas cuatro corridas hizo falta entre 1,35 y 2,22 veces esa cantidad. El capítulo registra también dos limitaciones. Los criterios de llegada (C1) y de cierre de la misión (C2) no son independientes en este tipo de fallo, porque el coordinador declara fallida la misión con el mismo umbral de 0,25~m. Y tres misiones se ejecutaron dos veces sin que quedara registrado el motivo.

\section{Aporte a los objetivos específicos}
\label{sec:aporte}

\subsection{Relación entre actividades y objetivos}

\begin{table}[H]
    \centering
    \small
    \caption{Trazabilidad entre las actividades de la semana y los objetivos específicos.}
    \label{tab:trazabilidad}
    \begin{tabular}{@{}p{4.6cm}cp{7.4cm}@{}}
        \toprule
        \textbf{Actividad} & \textbf{Objetivo} & \textbf{Aporte} \\ \midrule
        Cancelación de misiones (RF-29) & OE3 & Último requisito del objetivo 3; con él quedan verificados sus seis requisitos \\ \addlinespace
        Comunicación entre los dos vehículos (RF-15) & OE2 & Primer requisito verificado sobre los vehículos reales; cierra el riesgo R11 \\ \addlinespace
        Información de avance en los dos pisos & OE2 & Descarta, con mediciones, que el vehículo construya el mapa del pasillo real recorriéndolo \\ \addlinespace
        Respuesta de tracción del vehículo & OE2 & Reduce lo pendiente de RF-14 a una medición concreta \\ \addlinespace
        Cierre de R14 y reproducibilidad de los registros & OE4 & Los resultados de la campaña pueden regenerarse a partir de sus datos \\ \addlinespace
        Capítulo de resultados & OE4 & Identifica el modo de fallo de la campaña y lo cuantifica \\ \bottomrule
    \end{tabular}
\end{table}

\subsection{Estado de los objetivos específicos}

\begin{table}[H]
    \centering
    \small
    \caption{Avance por objetivo específico al cierre de la semana 23.}
    \label{tab:objetivos}
    \begin{tabular}{@{}cp{2.8cm}p{4.6cm}p{4.2cm}@{}}
        \toprule
        \textbf{Objetivo} & \textbf{Avance} & \textbf{Aporte de la semana} & \textbf{Pendiente} \\ \midrule
        OE1 & 100~\% & --- & --- \\ \addlinespace
        OE2 & 55~\% $\rightarrow$ 65~\% & RF-15 verificado; pasillos de los dos pisos medidos & La odometría sobre el vehículo real \\ \addlinespace
        OE3 & 90~\% $\rightarrow$ 100~\% & RF-29 verificado & --- \\ \addlinespace
        OE4 & 85~\% & Capítulo de resultados escrito; R14 cerrado & RF-27, la demostración física \\ \bottomrule
    \end{tabular}
\end{table}

El promedio simple del avance de los cuatro objetivos es 87,5~\%, frente al 71,9~\% que corresponde al calendario (semana 23 de 32). Están verificados 29 de los 36 requisitos.

Esta diferencia no representa un margen de tiempo. Todo lo pendiente depende de la plataforma física, y su avance no depende solo de las horas de trabajo.

\section{Estado del cronograma y trabajo previsto}
\label{sec:cronograma}

\subsection{Cumplimiento del criterio de cierre}

\begin{table}[H]
    \centering
    \small
    \caption{Criterio de cierre de la semana 23 y su evidencia.}
    \label{tab:cierre}
    \begin{tabular}{@{}p{5.4cm}cp{6.4cm}@{}}
        \toprule
        \textbf{Parte del criterio} & \textbf{Estado} & \textbf{Evidencia} \\ \midrule
        El sistema funciona de extremo a extremo en simulación & Cumplida & Corridas de RF-29 del 14 de septiembre; misión operada desde un teléfono en la semana 22 \\ \addlinespace
        Al menos una corrida física completa & No cumplida & La primera navegación sobre el vehículo real fue el 24 de septiembre, en la semana 24 \\ \addlinespace
        Repositorio etiquetado & Cumplida & Etiqueta \texttt{v0.4-implementacion-congelada}, 18 de septiembre \\ \bottomrule
    \end{tabular}
\end{table}

\subsection{Actividades planificadas y ejecutadas}

De las cuatro actividades planificadas se ejecutaron tres: la verificación y corrección de fallas, la congelación del código y la emisión del informe de la semana 22. El ensayo del protocolo con los vehículos reales no se hizo porque el vehículo aún no publicaba su odometría. Se realizaron además cuatro actividades no planificadas: la medición de la comunicación con el segundo vehículo, la medición de los pasillos de los dos pisos, la medición de la respuesta de tracción y el capítulo de resultados.

\subsection{Trabajo previsto para la semana 24}

Para la semana 24 se prevé decidir si se intenta la demostración física con los vehículos reales (decisión GO/NO-GO) y publicar la odometría en el vehículo, que es la ruta crítica. También se consolidará el conjunto de datos de la campaña y se emitirá este informe.

\section{Conclusiones}

\begin{enumerate}[leftmargin=*]
    \item El código quedó congelado el 18 de septiembre con los objetivos 1 y 3 completos. Lo pendiente de los objetivos 2 y 4 depende de la plataforma física.
    \item El requisito de comunicación entre vehículos (RF-15) se verificó el mismo día en que llegó el segundo vehículo: latencia mediana de 7,61~ms frente a un límite de 100~ms, sin pérdida de mensajes.
    \item El objetivo 3 quedó completo con la verificación de la cancelación de misiones (RF-29): el vehículo regresó a 0,121~m de la escalera, dentro del criterio de 0,25~m.
    \item La información de avance en los pasillos reales es de 5,1~\% en el piso 1 y 5,9~\% en el piso 2, por debajo del 6,8~\% del pasillo simulado cuyo mapa falló. Queda descartado que el vehículo construya el mapa del pasillo recorriéndolo, y la única ruta crítica es publicar su odometría.
    \item La corrida física completa que exigía el criterio de cierre no se realizó, y esa parte del criterio se declara no cumplida.
\end{enumerate}

\section{Anexos}

Las cifras de este informe proceden de los registros y documentos de evidencia del repositorio del proyecto, que se listan a continuación. El mismo informe está disponible en Markdown, en \texttt{Entregable\_semana\_23.md}, para leerlo directamente en el repositorio.

\begin{thebibliography}{99}

\bibitem{rf29}
S. Hernández Ávila y J. A. Mejía León,
``RF-29: corrida con cancelación y corrida de control,''
Repositorio del proyecto, \texttt{Documentos/Evidencia/registros/S23\_RF29\_cancelada.json} y \texttt{S23\_RF29\_control.json}, septiembre de 2026.

\bibitem{rf15}
S. Hernández Ávila y J. A. Mejía León,
``RF-15 medido entre los dos vehículos,''
Repositorio del proyecto, \texttt{Documentos/Evidencia/registros/S23\_RF15\_carroA\_carroB.json}, septiembre de 2026.

\bibitem{pisos}
S. Hernández Ávila,
``Información de avance del pasillo real, pisos 1 y 2,''
Repositorio del proyecto, \texttt{Documentos/Evidencia/S23\_informacion\_avance\_piso1.md} y \texttt{S23\_informacion\_avance\_piso2.md}, septiembre de 2026.

\bibitem{rf14}
S. Hernández Ávila,
``RF-14: respuesta de tracción medida sobre el vehículo,''
Repositorio del proyecto, \texttt{Documentos/Evidencia/S23\_campo\_traccion\_RF14.md}, septiembre de 2026.

\bibitem{mapa}
S. Hernández Ávila y J. A. Mejía León,
``Mapa del trabajo restante hasta la sustentación,''
Repositorio del proyecto, \texttt{Documentos/MAPA\_TRABAJO\_RESTANTE.md}, septiembre de 2026.

\bibitem{reproducibilidad}
S. Hernández Ávila,
``Reproducibilidad de los registros de la campaña,''
Repositorio del proyecto, \texttt{Documentos/Evidencia/S23\_reproducibilidad\_de\_los\_registros.md}, septiembre de 2026.

\bibitem{arranque}
S. Hernández Ávila,
``Guía de arranque del sistema,''
Repositorio del proyecto, \texttt{Documentos/GUIA\_ARRANQUE.md}, septiembre de 2026.

\bibitem{resultados}
S. Hernández Ávila y J. A. Mejía León,
``Resultados de la campaña de evaluación en simulación (OE4),''
Repositorio del proyecto, \texttt{Documentos/RESULTADOS\_OE4\_SIMULACION.md}, septiembre de 2026.

\end{thebibliography}

\end{document}
